\documentclass[../main.tex]{subfiles}
\graphicspath{{\subfix{../images/}}}
\begin{document}
	Un’interrogazione è una funzione $E(r)$ che applicata ad istanze $r$ di una base di dati produce una relazione su un dato insieme di attributi $X$.\\\\
	Le interrogazioni su uno schema di base di dati $R$ in algebra relazionale sono espressioni i cui atomi (le variabili) sono relazioni in $R$.
	\subsection{Esempio}
	\begin{center}
		\begin{tabular}{ |c|c|c|c| }
			\hline
			\multicolumn{4}{|c|}{Impiegati}\\
			\hline
			\textbf{Matricola} & \textbf{Nome} & \textbf{Età} & \textbf{Stipendio}\\
			\hline
			7309 & Rossi & 34 & 45\\
			5998 & Bianchi & 37 & 38\\
			9553 & Neri & 42 & 35\\
			5698 & Bruni & 43 & 42\\
			4076 & Mori & 45 & 50\\
			8123 & Lupi & 46 & 60\\
			\hline
		\end{tabular}
	\end{center}
	\begin{center}
		\begin{tabular}{ |c|c| }
			\hline
			\multicolumn{2}{|c|}{Supervisione}\\
			\hline
			\textbf{Impiegato} & \textbf{Capo} \\
			\hline
			7309 & 5698\\
			5998 & 5698\\
			9553 & 4076\\
			5698 & 4076\\
			4076 & 8123\\
			\hline
		\end{tabular}
	\end{center}
	Trovare matricola, nome, età e stipendio degli impiegati che guadagnano più di 40:
	\[SEL_{Stipendio > 40} (Impiegati)\]
	\begin{center}
		\begin{tabular}{ |c|c|c|c| }
			\hline
			\multicolumn{4}{|c|}{Impiegati}\\
			\hline
			\textbf{Matricola} & \textbf{Nome} & \textbf{Età} & \textbf{Stipendio}\\
			\hline
			7309 & Rossi & 34 & 45\\
			5698 & Bruni & 43 & 42\\
			4076 & Mori & 45 & 50\\
			8123 & Lupi & 46 & 60\\
			\hline
		\end{tabular}
	\end{center}
	Trovare matricola, nome ed età degli impiegati che guadagnano più di 40
	\[PROJ_{Matricola, Nome, \textit{Età}} (SEL_{Stipendio > 40} (Impiegati))\]
	\begin{center}
		\begin{tabular}{ |c|c|c| }
			\hline
			\multicolumn{3}{|c|}{Impiegati}\\
			\hline
			\textbf{Matricola} & \textbf{Nome} & \textbf{Età}\\
			\hline
			7309 & Rossi & 34\\
			5698 & Bruni & 43\\
			4076 & Mori & 45\\
			8123 & Lupi & 46\\
			\hline
		\end{tabular}
	\end{center}
	\subsection{Procedimento Logico}
	Trovare le matricole dei capi degli impiegati che guadagnano più di 40:
	\[PROJ_{Capo} ( SEL_{Stipendio>40} ( \text{Supervisione } JOIN_{Impiegato = Matricola} \text{ Impiegati} ) )\]
	Operando secondo lo schema che abbiamo visto:
	\begin{enumerate}
		\item Creo una singola tabella in cui ogni tupla contiene (fra gli altri) gli attributi Capo, Impiegato e Stipendio con un join fra Supervisione e Impiegati
		\item Seleziono le tuple con stipendio $>$ 40
		\item Proietto il risultato sull'attributo Capo
	\end{enumerate}
	\subsection{Algebra con valori nulli}
	Gli operatori logici vengono estesi con una logica a 3 valori: VERO (V), FALSO (F), SCONOSCIUTO (U)\\\\
	Tabelle di verità operatori logici:
	\begin{center}
		\begin{tabular}{ c|c }
			\textbf{NOT} & \\
			\hline
			\textbf{F} & V\\
			\textbf{U} & U\\
			\textbf{V} & F
		\end{tabular}
		\hfill
		\begin{tabular}{ c|c c c }
			\textbf{AND} & \textbf{V} & \textbf{U} & \textbf{F}\\
			\hline
			\textbf{V} & V & U & F \\
			\textbf{U} & U & U & F\\
			\textbf{F} & F & F & F
		\end{tabular}
		\hfill
		\begin{tabular}{ c|c c c }
			\textbf{OR} & \textbf{V} & \textbf{U} & \textbf{F}\\
			\hline
			\textbf{V} & V & V & V \\
			\textbf{U} & V & U & U\\
			\textbf{F} & V & U & F
		\end{tabular}
	\end{center}
	\subsection{Equivalenze di espressioni}
	Due espressioni sono equivalenti se:
	\begin{itemize}
		\item $E_1 \equiv_R E_2$ se $E_1(r)=E_2(r)$ per ogni istanza $r$ di $R$\\
		\textbf{Equivalenza dipendente dallo schema}
		\item $E_1\equiv E_2$ se $E_1 \equiv_R E_2$ per ogni schema $R$\\
		\textbf{Equivalenza assoluta}
	\end{itemize}
	L’equivalenza è importante in quanto consente di scegliere, a parità di risultato, l’operazione meno costosa.
	\begin{itemize}
		\item Atomizzazione delle selezioni\\
		$\sigma_{F_1 \wedge F_2} (E) \equiv \sigma_{F_1} (\sigma_{F_2} (E) )$
		\item Idempotenza delle proiezioni\\
		$\Pi_{X} (E) \equiv \Pi_{X} (\Pi_{XY} (E))$ qualunque sia $Y$
		\item Anticipazione della selezione rispetto al join\\
		$\sigma_F (E_1 \Diamond E_2) \equiv E_1 \Diamond (\sigma_F (E_2))$
		\item Anticipazione della proiezione rispetto al join:\\
		$\Pi_{X_1,Y_2} (E_1 \Diamond E_2) \equiv E_1 \Diamond \Pi_{Y_2} (E_2)$   {\color{red} NB: $Y_1 \subseteq X_1, Y_2 \subseteq X_2$ }\\
		($E_1$ definita su $X_1$, $E_2$ definita su $X_2$)\\
		(Vale solo se $Y_2$ contiene gli attributi coinvolti nella condizione di join)\\\\
		Allora (combinando con idempotenza delle proiezioni):\\
		$\Pi_Y (E_1 \Diamond_F E_2) \equiv \Pi_Y (\Pi_{Y_1} (E_1) \Diamond_F \Pi_{Y_2} (E_2))$\\
		Dove $Y_1$ e $Y_2$ contengono, rispettivamente, gli attributi di $X_1$ e $X_2$ compresi in $Y$ o coinvolti nella condizione $F$.\\\\
		In pratica è possibile ignorare, in ciascuna relazione, gli attributi non compresi in Y e non coinvolti nel join.
		\item Inglobamento di una selezione in un prodotto cartesiano a formare un join (theta-join):\\
		$\sigma_F (E_1 \Diamond E_2 ) \equiv E_1 \Diamond_F E_2$
		\item Tutti gli operatori binari eccetto la differenza godono delle proprietà associativa e commutativa.
		\item Distributività della selezione rispetto all'unione:\\
		$\sigma_F (E_1 \cup E_2) \equiv \sigma_F (E_1) \cup \sigma_F (E_2)$
		\item Distributività della selezione rispetto alla differenza:\\
		$\sigma_F (E_1 - E_2) \equiv \sigma_F (E_1) - \sigma_F (E_2)$
		\item Distributività della proiezione rispetto all'unione:\\
		$\Pi_{X} (E_1 \cup E_2) \equiv \Pi_{X} (E_1) \cup \Pi_{X} (E_2)$\\
		{ \color{red} NB La proiezione NON è distributiva rispetto alla differenza}
		\item Corrispondenze fra operatori insiemistici e selezioni complesse:\\
		$\sigma_{F_1 \vee F_2} (R) \equiv \sigma_{F_1} (R) \cup \sigma_{F_2} (R)$\\
		$\sigma_{F_1 \wedge F_2} (R) \equiv \sigma_{F_1} (R) \cap \sigma_{F_2} (R) \equiv \sigma_{F_1} (R) \Rightarrow \sigma_{F_2} (R)$\\
		$\sigma_{F_1 \wedge \neg F_2} (R) \equiv \sigma_{F_1} (R) - \sigma_{F_2} (R)$
		\item Proprietà distributiva del join rispetto all'unione:\\
		$E \Rightarrow (E_1 \cup E_2) \equiv (E \Rightarrow E_1) \cup (E \Rightarrow E_2)$
	\end{itemize}
	\subsection{Viste (Relazioni Derivate)}
	\begin{itemize}
		\item Rappresentazioni diverse per gli stessi dati (schema esterno):
		\begin{itemize}
			\item \textbf{Relazioni di base}: contenuto autonomo; fisicamente e originariamente contenute nella base di dati
			\item \textbf{Relazioni derivate}: relazioni il cui contenuto è funzione del contenuto di altre relazioni (definito per mezzo di interrogazioni)
			\item \textbf{Relazioni Virtuali (Viste)}: Relazioni definite mediante funzioni o espressioni del linguaggio di interrogazione, non realmente presenti nella base di dati ma utilizzabili come se lo fossero. \\
			Devono essere ricalcolate tutte le volte.
			\item \textbf{Viste materializzate}: Relazioni virtuali effettivamente inserite nella base di dati.\\
			Immediatamente disponibili ma critiche per il mantenimento dell’allineamento con le relazioni da cui derivano.\\
			Non sono supportate dai DBMS.
		\end{itemize}
	\end{itemize}
	\subsubsection{Vantaggi}
	\begin{itemize}
		\item Permettono di mostrare a un utente le sole componenti della base di dati che interessano
		\item Espressioni molto complesse possono essere definite come viste e rese molto più leggibili
		\item Sicurezza: è possibile definire diritti di accesso anche per una sola vista (quindi per una particolare porzione della base di dati)
		\item In caso di ristrutturazione della base di dati, le “vecchie” relazioni possono essere di nuovo ricavate mediante viste, consentendo l’uso di applicazioni che fanno riferimento al vecchio schema
	\end{itemize}
	\subsubsection{Esempio}
	\begin{center}
		\begin{tabular}{ |c|c| }
			\hline
			\multicolumn{2}{|c|}{Afferenza}\\
			\hline
			\textbf{Impiegato} & \textbf{Reparto} \\
			\hline
			Rossi & A\\
			Neri & B\\
			Bianchi & B\\
			\hline
		\end{tabular}
		\begin{tabular}{ |c|c| }
			\hline
			\multicolumn{2}{|c|}{Direzione}\\
			\hline
			\textbf{Reparto} & \textbf{Capo} \\
			\hline
			A & Mori\\
			B & Bruni\\
			\hline
		\end{tabular}
	\end{center}
	Una possibile vista:
	\[Supervisione(Impiegato, Capo) = PROJ_{Impiegato,Capo} (Afferenza \; JOIN \; Direzione)\]
	\subsubsection{Interrogazioni sulle viste}
	Sono eseguite sostituendo alla vista la sua definizione:
	\[SEL_{Capo = 'Leoni'} (Supervisione)\]
	viene eseguita come
	\[SEL_{Capo = 'Leoni'} (PROJ_{Impiegato, Capo} (Afferenza \; JOIN \; Direzione))\]
	\subsubsection{Selezione con condizione valutata su tuple diverse della stessa relazione}
	Trovare gli impiegati con lo stesso nome.\\\\
	Data la relazione:
	\[Impiegato (Nome, Cognome, MatricolaImpiegato)\]
	Definisco una vista:
	\[Impiegato1 (Nome1, Cognome1, MatricolaImpiegato1)\] 
	\[=\] 
	\[REN_{Nome1, Cognome1, MatricolaImpiegato1 \leftarrow Nome, Cognome, MatricolaImpiegato} (Impiegato)\]
	Interrogazione che ottiene il risultato voluto:
	\[PROJ_{Nome, Cognome} (SEL_{Nome = Nome1} (Impiegato \; JOIN \; Impiegato1))\]
	\begin{center}
		\textbf{Prodotto Cartesiano}
	\end{center}
	\subsubsection{Viste come strumenti di programmazione}
	Trovare gli impiegati che hanno lo stesso capo di Rossi
	\begin{itemize}
		\item Senza vista: (I "..." significano che va a capo)\\
		$PROJ_{Impiegato} ((Afferenza \; JOIN \; Direzione) \; JOIN \; REN_{ImpR, RepR \leftarrow Impiegato, Reparto} ...\\... (SEL_{Impiegato = 'Rossi'} (Afferenza \; JOIN \; Direzione)))$
		\item Con Vista:\\
		$PROJ_{Impiegato} ((Supervisione) \; JOIN \; REN_{ImpR, RepR \leftarrow Impiegato, Reparto} (SEL_{Impiegato = 'Rossi'} (Supervisione)))$
	\end{itemize}
	\textbf{NB} In questo esempio, rispetto al precedente:
	\[Supervisione(Impiegato,Reparto,Capo)= Afferenza \; JOIN \;  Direzione\]
	È senza proiezione, quindi contiene anche l’attributo Reparto!
	\subsubsection{Viste ed Aggiornamenti}
	\begin{center}
		\begin{tabular}{ |c|c| }
			\hline
			\multicolumn{2}{|c|}{Afferenza}\\
			\hline
			\textbf{Impiegato} & \textbf{Reparto} \\
			\hline
			Rossi & A\\
			Neri & B\\
			Verdi & A\\
			\hline
		\end{tabular}
		\begin{tabular}{ |c|c| }
			\hline
			\multicolumn{2}{|c|}{Direzione}\\
			\hline
			\textbf{Reparto} & \textbf{Capo} \\
			\hline
			A & Mori\\
			B & Bruni\\
			C & Bruni\\
			\hline
		\end{tabular}
	\end{center}
	\begin{center}
		\begin{tabular}{ |c|c| }
			\hline
			\multicolumn{2}{|c|}{Supervisione}\\
			\hline
			\textbf{Impiegato} & \textbf{Capo} \\
			\hline
			Rossi & Mori\\
			Neri & Bruni\\
			Verdi & Mori\\
			\hline
		\end{tabular}
	\end{center}
	Vogliamo inserire, nella vista, il fatto che Lupi ha come capo Bruni; oppure che Belli ha come capo Falchi; come facciamo?
	\begin{itemize}
		\item "Aggiornare una vista": modificare le relazioni di base in modo che la vista, "ricalcolata", rispecchi l'aggiornamento
		\item L'aggiornamento sulle relazioni di base corrispondente a quello sulla vista deve essere univoco
		\item In generale però non è univoco!
		\item Sulle viste sono ammissibili pochissimi aggiornamenti
	\end{itemize}
\end{document}